% !TEX program = xelatex
% !BIB program = biber
\documentclass[pl,12pt]{TAGH/t_agh_av}
\usepackage{float}
\usepackage[hidelinks]{hyperref}
\addbibresource{bibliography.bib}

\author{Bart\l{}omiej Pomiet\l{}o}
\shortauthor{B.P.}

\titlePL{Analiza narz\k{e}dzi dla przeciwdzia\l{}ania dezinformacji w obszarze multimedi\'{o}w}
\titleEN{Analysis of tools for countering disinformation in the field of multimedia}

\shorttitlePL{Analiza narz\k{e}dzi przeciwdzia\l{}ania dezinformacji w multimediach }
\shorttitleEN{Tools for countering multimedia disinformation }
\shortauthor{B.~Pomiet\l{}o }

\thesistype{Praca Magisterska}
\supervisor{dr~in\.{z}. Jan Derkacz}
\degreeprogramme{Cyberbezpiecze\'{n}stwo}

\date{2026}
\department{Katedra Telekomunikacji}
\faculty{Wydzia\l{} Informatyki, Elektroniki i Telekomunikacji}
\acknowledgements{}

\begin{document}
\titlepages
\selectlanguage{polish}

% -------------------------------------------------------
% ABSTRAKTY
% -------------------------------------------------------
\begin{abstract}
% TODO: napisac streszczenie angielskie
\phantom{x}
\end{abstract}

\begin{abstract}
% TODO: napisac streszczenie polskie
\phantom{x}
\end{abstract}

% -------------------------------------------------------
% SPISY
% -------------------------------------------------------
\tableofcontents
\clearpage

\chapter*{Wykaz tabel}
\addcontentsline{toc}{chapter}{Wykaz tabel}
\renewcommand{\listtablename}{}
\listoftables

\chapter*{Wykaz rysunk\'{o}w}
\addcontentsline{toc}{chapter}{Wykaz rysunk\'{o}w}
\renewcommand{\listfigurename}{}
\listoffigures

% -------------------------------------------------------
% WYKAZ SKROTOW
% -------------------------------------------------------
\chapter*{Wykaz skr\'{o}t\'{o}w}
\addcontentsline{toc}{chapter}{Wykaz skr\'{o}t\'{o}w}
\begin{tabular}{ll}
\textbf{AI}    & Artificial Intelligence (Sztuczna inteligencja) \\
\textbf{AUC}   & Area Under the Curve \\
\textbf{CNN}   & Convolutional Neural Network \\
\textbf{DNN}   & Deep Neural Network \\
\textbf{FFT}   & Fast Fourier Transform \\
\textbf{FP}    & False Positive \\
\textbf{FPR}   & False Positive Rate \\
\textbf{FN}    & False Negative \\
\textbf{GAN}   & Generative Adversarial Network \\
\textbf{LLM}   & Large Language Model \\
\textbf{ML}    & Machine Learning \\
\textbf{OF}    & Optical Flow \\
\textbf{TP}    & True Positive \\
\textbf{TPR}   & True Positive Rate (Recall) \\
\textbf{UE}    & Unia Europejska \\
\textbf{ViT}   & Vision Transformer \\
\end{tabular}

% -------------------------------------------------------
% ROZDZIALY
% -------------------------------------------------------
\include{chapters/01_wprowadzenie}
\include{chapters/02_tlo_literaturowe}
\include{chapters/03_zagrozenia_i_techniki}
\chapter{Narz\k{e}dzia i~metody detekcji}
\label{ch:narzedzia}

\section{Zbi\'{o}r danych testowych}
\label{sec:zbior_danych}

Do ewaluacji narz\k{e}dzi detekcji tre\'{s}ci generowanych przez sztuczn\k{a}
inteligencj\k{e} przygotowano zbi\'{o}r 22~film\'{o}w wideo, obejmuj\k{a}cy
materia\l{}y syntetyczne, edytowane oraz autentyczne nagrania stanowi\k{a}ce
klasy negatywn\k{a} (punkt odniesienia). Tabela~\ref{tab:dataset} prezentuje
pe\l{}ne zestawienie plik\'{o}w.

\begin{table}[H]
\centering
\caption{Zbi\'{o}r testowy film\'{o}w wideo}
\label{tab:dataset}
\small
\begin{tabular}{|l|r|l|l|}
\hline
\textbf{Nazwa pliku} & \textbf{Rozmiar} & \textbf{Typ} & \textbf{Model/\'{z}r\'{o}d\l{}o} \\
\hline
\multicolumn{4}{|c|}{\textit{Runway Gen-4.5 (wideo AI)}} \\
\hline
runway\_gen45\_clip\_02\_runway\_h264.mp4 & 8,1 MB & Syntetyczny & Runway Gen-4.5 \\
runway\_gen45\_clip\_03\_runway\_h264.mp4 & 8,1 MB & Syntetyczny & Runway Gen-4.5 \\
runway\_gen45\_clip\_04\_runway\_h264.mp4 & 6,4 MB & Syntetyczny & Runway Gen-4.5 \\
runway\_gen45\_clip\_05\_runway\_h264.mp4 & 2,2 MB & Syntetyczny & Runway Gen-4.5 \\
\hline
\multicolumn{4}{|c|}{\textit{OpenAI Sora (wideo AI)}} \\
\hline
kobe\_tuition\_90k\_sora.mp4 & 326 KB & Syntetyczny & OpenAI Sora \\
not\_let\_off\_plane\_sora.mp4 & 300 KB & Syntetyczny & OpenAI Sora \\
doctor\_ankle\_sora.mp4 & 296 KB & Syntetyczny & OpenAI Sora \\
1PaoWKvcJP0\_Sound\_on\_for\_Sora\_2.mp4 & 12,7 MB & Syntetyczny & OpenAI Sora \\
\hline
\multicolumn{4}{|c|}{\textit{Runway Gen-3/Gen-4 Turbo (wideo AI)}} \\
\hline
Gen-4\_Turbo\_candle\_flame.mp4 & 1,2 MB & Syntetyczny & Runway Gen-4 Turbo \\
Abstract\_fluid\_simulation\_neon.mp4 & 780 KB & Syntetyczny & Runway Gen-3 \\
AbEQv\_Water\_4k\_AI\_short\_film\_Gen3.mp4 & 10,3 MB & Syntetyczny & Runway Gen-3 \\
AwKSrJFvdps\_Introducing\_Gen45.mp4 & 14,2 MB & Syntetyczny & Runway Gen-4.5 \\
A\_dragon\_flying\_fantasy\_castle.mp4 & 481 KB & Syntetyczny & Runway Gen-3 \\
A\_person\_running\_rainy\_street.mp4 & 1,1 MB & Syntetyczny & Runway Gen-3 \\
\hline
\multicolumn{4}{|c|}{\textit{Inne modele AI / materia\l{}y promocyjne}} \\
\hline
y2dynet9Gvg\_Text\_to\_Video\_Demo.mp4 & 3,6 MB & Syntetyczny & r\'{o}\.{z}ne modele \\
Zb3tffmBPRE\_Luma\_Dream\_Machine.mp4 & 6,1 MB & Syntetyczny & Luma Dream Machine \\
\hline
\multicolumn{4}{|c|}{\textit{Materia\l{}y zmanipulowane i~deepfake}} \\
\hline
000000biden\_fake.mp4 & 3,6 MB & Deepfake & face-swap \\
000000kot\_fake.mp4 & 2,7 MB & Zmanipulowany & -- \\
000000passport\_fake.mp4 & 2,2 MB & Zmanipulowany & -- \\
000000pilot\_fake.mp4 & 2,3 MB & Zmanipulowany & -- \\
\hline
\multicolumn{4}{|c|}{\textit{Materia\l{}y autentyczne (klasa negatywna)}} \\
\hline
000000biden\_real.mp4 & 3,3 MB & Autentyczny & nagranie kamerowe \\
\hline
\multicolumn{4}{|c|}{\textit{Materia\l{}y kontrolne (t\l{}o sieciowe)}} \\
\hline
\_gxXv5qlEdw\_BBC\_News\_Blooper.mp4 & 4,9 MB & Autentyczny & BBC News \\
\_qi\_p3VXzCQ\_Lower\_Third\_Motion.mp4 & 3,4 MB & Autentyczny & motion graphics \\
\hline
\end{tabular}
\end{table}

Zbi\'{o}r sk\l{}ada si\k{e} z~\textbf{24~plik\'{o}w} (\l{}\k{a}cznie ok.\ 90~MB), w~tym:
\begin{itemize}
  \item \textbf{16 film\'{o}w w~pe\l{}ni generowanych przez AI} (Runway, Sora, Luma),
  \item \textbf{4 materia\l{}y zmanipulowane lub deepfake},
  \item \textbf{2 autentyczne nagrania kamerowe} (klasa negatywna),
  \item \textbf{2 autentyczne materia\l{}y kontrolne} (motion graphics, nagranie BBC).
\end{itemize}

R\'{o}\.{z}norodno\'{s}\'{c} model\'{o}w generatywnych (Runway Gen-3, Gen-4, Gen-4.5,
OpenAI Sora, Luma Dream Machine) oraz zr\'{o}\.{z}nicowanie rozmiar\'{o}w plik\'{o}w
(od 296~KB do 14,2~MB) pozwala oceni\'{c} odporno\'{s}\'{c} narz\k{e}dzi na~r\'{o}\.{z}ne
warunki kompresji i~styl generowania.

% -------------------------------------------------------------------
\section{Przegl\k{a}d istniej\k{a}cych narz\k{e}dzi internetowych}
\label{sec:przeglad_narzedzi}

Na rynku dost\k{e}pnych jest kilka publicznie dost\k{e}pnych serwis\'{o}w internetowych
pretenduj\k{a}cych do detekcji tre\'{s}ci generowanych przez AI lub zmanipulowanych
multimedi\'{o}w. W~niniejszej sekcji om\'{o}wiono narz\k{e}dzia testowane w~ramach
niniejszej pracy: DetectVideo.ai oraz Undetectable.ai (wyniki omawiane w~kolejnej
sekcji). Analiza obejmuje podej\'{s}cie techniczne, dost\k{e}pno\'{s}\'{c}, ograniczenia
oraz wyniki na~zbiorze testowym.

% -------------------------------------------------------------------
\subsection{DetectVideo.ai}
\label{sec:detectvideo}

DetectVideo.ai to serwis internetowy oferuj\k{a}cy detekcj\k{e} film\'{o}w wideo pod k\k{a}tem
pochodzenia z~modeli generatywnych AI. Narz\k{e}dzie zwraca wynik
w~postaci procentowego prawdopodobie\'{n}stwa tego, \.{z}e przes\l{}any plik
zosta\l{} wygenerowany lub zmanipulowany przez AI.

\subsubsection{Wyniki na zbiorze testowym}

Przeprowadzono test wszystkich 24~plik\'{o}w ze~zbioru opisanego
w~sekcji~\ref{sec:zbior_danych}. Tabela~\ref{tab:detectvideo_results}
przedstawia szczeg\'{o}\l{}owe wyniki.

\begin{table}[H]
\centering
\caption{Wyniki DetectVideo.ai dla zbioru testowego}
\label{tab:detectvideo_results}
\small
\begin{tabular}{|l|r|c|}
\hline
\textbf{Nazwa pliku} & \textbf{Rozmiar} & \textbf{Wynik AI likelihood} \\
\hline
\multicolumn{3}{|c|}{\textit{Runway Gen-4.5}} \\
\hline
runway\_gen45\_clip\_05 & 2,2 MB & 55\% \\
runway\_gen45\_clip\_04 & 6,4 MB & 55\% \\
runway\_gen45\_clip\_02 & 8,1 MB & 55\% \\
runway\_gen45\_clip\_03 & 8,1 MB & 55\% \\
\hline
\multicolumn{3}{|c|}{\textit{OpenAI Sora}} \\
\hline
kobe\_tuition\_90k\_sora & 326 KB & 55\% \\
not\_let\_off\_plane\_sora & 300 KB & 55\% \\
doctor\_ankle\_sora & 296 KB & 55\% \\
1PaoWKvcJP0\_Sora & 12,7 MB & 55\% \\
\hline
\multicolumn{3}{|c|}{\textit{Runway Gen-4 Turbo / Gen-3}} \\
\hline
Gen-4\_Turbo\_candle\_flame & 1,2 MB & 55\% \\
Abstract\_fluid\_simulation\_neon & 780 KB & \textbf{72\%} \\
AbEQv\_Water\_4k\_AI\_short\_film\_Gen3 & 10,3 MB & \textbf{72\%} \\
A\_dragon\_flying\_fantasy\_castle & 481 KB & \textbf{72\%} \\
A\_person\_running\_rainy\_street & 1,1 MB & \textbf{72\%} \\
AwKSrJFvdps\_Gen45\_promo & 14,2 MB & 55\% \\
\hline
\multicolumn{3}{|c|}{\textit{Inne modele AI}} \\
\hline
y2dynet9Gvg\_Text\_to\_Video\_Demo & 3,6 MB & 55\% \\
Zb3tffmBPRE\_Luma\_Dream\_Machine & 6,1 MB & 55\% \\
\hline
\multicolumn{3}{|c|}{\textit{Deepfake i zmanipulowane}} \\
\hline
000000biden\_fake & 3,6 MB & 55\% \\
000000kot\_fake & 2,7 MB & 55\% \\
000000passport\_fake & 2,2 MB & \textbf{72\%} \\
000000pilot\_fake & 2,3 MB & 55\% \\
\hline
\multicolumn{3}{|c|}{\textit{Autentyczne (klasa negatywna)}} \\
\hline
000000biden\_real & 3,3 MB & 55\% \\
\_gxXv5qlEdw\_BBC\_News\_Blooper & 4,9 MB & 55\% \\
\_qi\_p3VXzCQ\_Lower\_Third\_Motion & 3,4 MB & 55\% \\
\hline
\end{tabular}
\end{table}

\subsubsection{Analiza wynik\'{o}w}

Wyniki przedstawione w~tabeli~\ref{tab:detectvideo_results} ujawniaj\k{a}
istotne ograniczenia narz\k{e}dzia. Serwis DetectVideo.ai zwr\'{o}ci\l{} wynik
\textbf{55\%} (,,AI likelihood'') dla \textbf{20 z~24~plik\'{o}w} (83\%),
w~tym r\'{o}wnie\.{z} dla pliku \texttt{000000biden\_real.mp4} --- autentycznego
nagrania kamerowego. Oznacza to, \.{z}e wynik 55\% stanowi~de~facto domy\'{s}lny
odpowied\'{z} narz\k{e}dzia w~przypadku braku pewno\'{s}ci --- jest to~wynik
,,nie wiem'', a~nie faktyczna ocena.

Wy\.{z}szy wynik (72\%) zosta\l{} przyznany jedynie 5~filmom:
\begin{itemize}
  \item materia\l{}om z~oczywistymi cechami generatywnymi
        (abstrakcyjna animacja p\l{}yn\'{o}w, smok lataj\k{a}cy nad zamkiem,
        efektowne AI city scene);
  \item filmowi promocyjnemu Runway Gen-3 (\textit{Water});
  \item sfalsyfikowanemu dokumentowi (\texttt{passport\_fake}) ---
        wskazanemu jako manipulowany dzi\k{e}ki analizie metadanych,
        a~nie tre\'{s}ci wizualnej.
\end{itemize}

Diagnoza jest jednoznaczna: \textbf{DetectVideo.ai jest praktycznie bezużyteczny
jako detektor film\'{o}w AI w~przypadku trudnych, nowych modeli generatywnych
(Runway Gen-4.5, OpenAI Sora)}. Narz\k{e}dzie nie potrafi odr\'{o}\.{z}ni\'{c}
autentycznego nagrania od syntetycznego w~80\% przypadk\'{o}w i~zwraca
domy\'{s}lny wynik ,,inconclusive'' zamiast podj\k{a}\'{c} decyzj\k{e}. Dla
niniejszego detektora fakt ten stanowi wa\.{z}ny punkt odniesienia~---
proponowane rozwi\k{a}zanie podejmuje jednoznaczn\k{a} decyzj\k{e} binarny dla
ka\.{z}dego pliku.

% Tutaj mozna wstawic zrzut ekranu interfejsu DetectVideo.ai
% \begin{figure}[H]
%   \centering
%   \includegraphics[width=0.85\textwidth]{figures/detectvideo_screenshot.png}
%   \caption{Interfejs DetectVideo.ai z wynikiem dla pliku 000000passport\_fake.mp4 (72\% AI likelihood)}
%   \label{fig:detectvideo_screenshot}
% \end{figure}

% -------------------------------------------------------------------
\subsection{Undetectable.ai}
\label{sec:undetectable}

% Sekcja zostanie uzupelniona po otrzymaniu wynikow z Undetectable.ai
\textit{Wyniki dla narz\k{e}dzia Undetectable.ai zostan\k{a} om\'{o}wione
po zakończeniu test\'{o}w na zbiorze danych.}

% -------------------------------------------------------------------
\section{Por\'{o}wnanie narz\k{e}dzi}
\label{sec:porownanie}

\subsection{Zalety i~wady analizowanych narz\k{e}dzi internetowych}
\label{sec:zalety_wady_internet}

Na podstawie przeprowadzonych test\'{o}w sformu\l{}owano nast\k{e}puj\k{a}c\k{a}
charakterystyk\k{e} dost\k{e}pnych narz\k{e}dzi komercyjnych:

\begin{table}[H]
\centering
\caption{Por\'{o}wnanie komercyjnych narz\k{e}dzi detekcji}
\label{tab:porownanie_narzedzi}
\small
\begin{tabular}{|p{3.5cm}|p{5.5cm}|p{5.5cm}|}
\hline
\textbf{Narz\k{e}dzie} & \textbf{Zalety} & \textbf{Wady} \\
\hline
DetectVideo.ai &
Prosty w~u\.{z}yciu, brak rejestracji, obs\l{}uguje popularne formaty wideo,
mo\.{z}liwe wskazanie sygna\l{}u metadanych &
80\% wynik\'{o}w to inconclusive 55\%; nie r\'{o}\.{z}nicuje nowych modeli
(Sora, Gen-4.5); brak wyja\'{s}nienia decyzji (XAI); faz wynik daje si\k{e}
zinterpretowa\'{c} jedynie jako ,,nie wiem'' \\
\hline
Undetectable.ai &
Znany z~detekcji tekstu AI; popularny w\'{s}r\'{o}d u\.{z}ytkownik\'{o}w &
G\l{}\'ownie przeznaczony do~tekstu; wsparcie wideo ograniczone;
wyniki do uzupe\l{}nienia \\
\hline
\end{tabular}
\end{table}

\subsection{Zalety i~wady proponowanego detektora}
\label{sec:zalety_wady_wlasny}

Proponowane w~niniejszej pracy podej\'{s}cie opiera si\k{e} na~fuzji
wielokana\l{}owej: analizy spektralnej (FFT), przep\l{}ywu optycznego
oraz metadanych. Poni\.{z}ej zestawiono jego g\l{}\'{o}wne w\l{}a\'{s}ciwo\'{s}ci.

\paragraph{Zalety}
\begin{itemize}
  \item \textbf{Decyzja binarna} --- detektor zawsze podejmuje jednoznaczn\k{a}
        decyzj\k{e} (AI / nie AI), eliminuj\k{a}c nierozstrzygni\k{e}te wyniki.
  \item \textbf{Wielosygna\l{}owo\'{s}\'{c}} --- po\l{}\k{a}czenie cech spektralnych,
        przep\l{}ywu optycznego i~metadanych zwi\k{e}ksza odporno\'{s}\'{c}
        na~manipulacje pojedynczym kana\l{}em.
  \item \textbf{Interpretowalny wynik} --- ka\.{z}da decyzja zawiera wag\k{e}
        sk\l{}adowych sygna\l{}\'{o}w, co umo\.{z}liwia analiz\k{e} przyczynow\k{a}.
  \item \textbf{Brak zale\.{z}no\'{s}ci od chmury} --- mo\.{z}liwo\'{s}\'{c} uruchomienia
        lokalnie, co istotne w~kontekstach wymagaj\k{a}cych prywatno\'{s}ci danych.
  \item \textbf{Skalowalny na nowe modele} --- po przetrenowaniu na~nowych danych
        mo\.{z}na dostosowa\'{c} do kolejnych generacji modeli AI.
\end{itemize}

\paragraph{Wady i~ograniczenia}
\begin{itemize}
  \item \textbf{Czuło\'{s}\'{c} na kompresj\k{e}} --- silna kompresja H.264/H.265
        mo\.{z}e wyt\l{}umi\'{c} sygna\l{}y spektralne.
  \item \textbf{Kr\'{o}tkie klipy} --- pliki poni\.{z}ej 5~sekund dostarczaj\k{a}
        niewystarczaj\k{a}cej liczby klatek do rzetelnej analizy przep\l{}ywu
        optycznego.
  \item \textbf{Zasoby obliczeniowe} --- analiza wielokana\l{}owa wymaga
        wi\k{e}kszej mocy obliczeniowej ni\.{z} proste narz\k{e}dzia oparte
        na~metadanych.
  \item \textbf{Ograniczenia generalizacji} --- narz\k{e}dzie mo\.{z}e gorzej
        dzia\l{}a\'{c} na~modelach generatywnych nieobj\k{e}tych zbiorem treningowym.
\end{itemize}

\include{chapters/05_implementacja_i_ewaluacja}
\chapter{Wnioski i kierunki dalszych badań}
\label{chap:wnioski}

\section{Podsumowanie wyników}

Celem niniejszej pracy była analiza narzędzi i metod wykrywania treści generowanych
przez sztuczną inteligencję w obszarze multimediów wideo, ze szczególnym uwzględnieniem
praktycznej skuteczności dostępnych detektorów oraz zaproponowanie podejścia
wielosygnałowego jako alternatywy dla rozwiązań komercyjnych.

Przeprowadzone eksperymenty wykazały, że popularne narzędzie DetectVideo.ai zwraca
wynik 55\% (inconclu\-sive) dla ponad 80\% testowanych materiałów, w tym dla nagrań
autentycznych, co czyni je praktycznie bezużytecznym do binarnej klasyfikacji treści.
Jedynie filmy o wyraźnych artefaktach wizualnych typowych dla generatywnych
modeli (np.\ abstrakcyjne symulacje płynów, sceny fantasy) uzyskały wynik 72\%,
co nadal nie przekracza przyjętego w literaturze progu pewności detekcji.
Proponowane podejście wielosygnałowe --- łączące analizę spektralną, przepływ
optyczny oraz analizę metadanych --- wykazuje zdolność do podjęcia jednoznacznej
decyzji binarnej dla każdego badanego pliku, co stanowi jego kluczową przewagę
praktyczną.

\section{Kierunki dalszych badań}

Niniejsza praca stanowi punkt wyjścia do szeroko zakrojonych badań w obszarze
detekcji multimediów generowanych przez AI. Poniżej przedstawiono kierunki,
które zdaniem autora posiadają największy potencjał naukowy i praktyczny.

\subsection{Rozszerzenie zbioru danych testowych}

Podstawowym ograniczeniem przeprowadzonych eksperymentów jest stosunkowo niewielki
zbiór 24 plików wideo. Dalsze badania powinny obejmować:

\begin{itemize}
  \item \textbf{Większe zbiory darmowe} --- takie jak FaceForensics++~\cite{Rossler2019FaceForensics},
        DFDC (Deepfake Detection Challenge Dataset) czy WildDeepfake, liczące dziesiątki
        tysięcy klipów z różnymi technikami syntezy;
  \item \textbf{Płatne bazy danych} --- komercyjne zbiory, np.\ oferowane przez Veritone
        lub Nielsen, zawierające materiały z kontrolowanymi metadanymi produkcji,
        co umożliwia precyzyjną weryfikację etykiet;
  \item \textbf{Materiały z platform społecznościowych} --- filmy pobrane z TikToka,
        YouTube Shorts i Instagrama, które przeszły przez agresywną rekompresję
        transkodowania platformy, stanowiącą jedno z najtrudniejszych wyzwań
        dla detektorów sygnałowych;
  \item \textbf{Zróżnicowanie modeli generatywnych} --- włączenie materiałów z nowszych
        modeli (Sora 2, Runway Gen-5, Kling, Pika 2.0) w miarę ich dostępności,
        aby śledzić ewolucję artefaktów.
\end{itemize}

\subsection{Ewaluacja płatnych narzędzi komercyjnych}

W ramach niniejszej pracy przetestowano wyłącznie narzędzia dostępne bezpłatnie.
Pełna analiza porównawcza powinna objąć płatne rozwiązania klasy enterprise, takie jak:
SynthID Detector (Google DeepMind), Hive Moderation, Reality Defender oraz Intel
FakeCatcher. Szczególnie interesujące byłoby zbadanie, czy wyższy koszt subskrypcji
przekłada się na rzeczywistą poprawę czułości w stosunku do proponowanego
podejścia wielosygnałowego.

\subsection{Odporność na ataki adwersarialne}

Praktyczny detektor musi być odporny na celowe działania twórców deepfake'ów.
Kierunkiem wartym zbadania jest testowanie detektora wobec technik adversarial,
takich jak dodawanie szumu gaussowskiego przed publikacją, kompresja
H.265 z niskim bitrate oraz preprocessing przez filtry stylistyczne
(np.\ efekt ,,film grain'').
Badania w tym obszarze pozwolą na identyfikację najsłabszych ogniw proponowanego
potoku i ukierunkowanie dalszych prac na ich wzmocnienie.

\subsection{Metody uczenia maszynowego}

Obecna implementacja opiera się na metodach sygnałowych opartych na
regułach heurystycznych. Naturalnym kierunkiem rozwoju jest zastąpienie
lub uzupełnienie modułu decyzyjnego klasyfikatorem uczonym nadzorowanym ---
np.\ lekką siecią CNN na cechach spektralnych lub modelem XGBoost
na wektorze cech numerycznych --- co potencjalnie zwiększy zdolność generalizacji
na nowych modelach generatywnych bez konieczności ręcznej aktualizacji progów.

\subsection{Detekcja w czasie rzeczywistym i integracja z platformami}

Obecna architektura przetwarza pliki wsadowo. Kolejnym krokiem praktycznym
byłoby opracowanie wersji strumieniowej zdolnej do analizy wideo na żywo
(np.\ podczas transmisji na platformach streamingowych) lub wdrożonej jako
wtyczka do CMS redakcyjnych, co odpowiada na rosnące zapotrzebowanie
branży mediowej na narzędzia weryfikacji treści w czasie rzeczywistym.

\subsection{Cyfrowe znaki wodne jako proaktywna ochrona}

Obok reaktywnej detekcji warto rozwinąć badania nad proaktywnym
oznaczaniem treści AI na poziomie modelu generatywnego (C2PA,
SynthID). Analiza skuteczności przetrwania tych znaków po typowych
operacjach edycyjnych (przycięcie, zmiana rozdzielczości, transkodowanie)
stanowi otwarty problem badawczy z bezpośrednim przełożeniem na regulacje prawne UE
dotyczące oznaczania treści syntetycznych~\cite{EUAIAct2024}.


% -------------------------------------------------------
% BIBLIOGRAFIA
% -------------------------------------------------------
\printbibliography[heading=bibintoc, title={Bibliografia}]

\end{document}
